\section{Introduction}
A renewal service provider may observe many histories while its terminal gateway accepts only a small menu of commands. The provider must choose that menu and allocate an accepted service obligation across histories. Participation limits expected service on each history, a separate ceiling limits every realized delivery, and installing a command can incur a fixed charge. These constraints make the menu and the promised service targets interdependent. Optimizing commands for prescribed targets does not solve the joint design problem.

We develop a joint-design method that preserves the original command budget, charges, aggregate obligation, and individual constraints. Its first component is an exact price oracle on a box of target intervals. A mandatory lower target can make commands after the price-optimal peak indispensable. An increasing prefix accounts for upper-target crossings, a decreasing suffix accounts for lower-target crossings, and their common peak joins two path recurrences. The oracle gives valid upper bounds over every menu in the box, not just over a menu supplied by a heuristic. Adaptive partitioning then yields an interval containing the global joint optimum, with finite termination at every positive additive tolerance.

The second component reduces the history dimension without requiring identical histories. We group histories that admit the same candidate commands but can have different participation caps and intermediate-service costs. Weighted average caps preserve aggregate capacity; minimum cost curvatures yield an optimistic reduced model. One original minimum-cap history guards against selecting a menu whose lowest command violates an individual cap. A constructive recovery procedure clips reduced targets at the original caps and redistributes the deficit within each group. It preserves the accepted obligation exactly and returns lotteries satisfying every original realization ceiling. We bound the loss from cap dispersion and cost dispersion separately, and combine it with the reduced solver's global interval.

This construction provides a variable-history additive approximation scheme for the normalized quadratic model when eligibility complexity and primitive scales are bounded. The number of groups needed for a prescribed error depends on accuracy and on the number of distinct eligible command sets, not on the number of histories. With a common realization ceiling, the eligibility count stays fixed as both the history population and catalog grow. The worst-case accuracy dependence remains large. Coarser practical partitions are therefore judged by their measured original-space certificate, and can be refined when that interval is too wide. A reduced solver's completion flag alone never certifies the original problem.

An independent checker verifies the group construction, recovery identities, original lotteries, price-path inequalities, and the complete target cover. Our numerical study includes distinct-history populations rather than only exact replicas. It compares certified coarsening with unaggregated adaptive search, the inherited priced-prefix method, exhaustive menu enumeration, and an uneliminated mixed-integer formulation. All declared cases and unresolved intervals are retained. Separate stress instances show that an exact reduced solution can still leave a substantial original-space gap, and that refinement removes this gap. These synthetic experiments evaluate algorithms and certificates, not an estimated operational population.

\paragraph{Relation to established methods.}\label{sec:r37-position}
Classical interval quantization and matrix searching underpin the saturated and conditional path algorithms \citep{Wu1991,NielsenNock2014,GronlundEtAl2017}. Pricing, constrained paths, and branch-and-bound are established tools \citep{HandlerZang1980,LawlerWood1966,Lemarechal2001}. Spatial branch-and-bound also exploits locally tight relaxations of separable piecewise-linear problems \citep{HubnerEtAl2025}. We do not claim those general mechanisms as new. Our two-sided oracle handles a shared, charged, cardinality-constrained menu under heterogeneous eligibility, with an anchor condition that ensures local tightness.

Aggregation error analysis is likewise established. \citet{LiBertsekas2025} bound value-function error in aggregated discounted dynamic programming, and \citet{Fattahi2025} studies approximation and error components in jointly designing charging programs and managing a heterogeneous participant population. Our result concerns a different coupling: one installed menu and an exact root obligation must survive aggregation and recovery while every original cap and realization ceiling remains enforceable. The operative ingredients are eligibility-preserving groups, the minimum-cap guard, the optimistic representative, and the explicit group-mass-preserving lift. The associated complexity guarantee does not assume a small number of exactly repeated types.

The broader contractual and continuous-design contributions remain part of the paper. Canonical allocation connects expected and pathwise participation; continuous quadratic design, Monge acceleration, conditional compression, full target nets, and charge-aware recovery retain their original assumptions and proofs. Table~\ref{tab:r46-contributions} separates these antecedents from two-sided joint decomposition and certified response coarsening. The article presents the joint-design argument first; the companion retains the full preceding mathematical chain and the preceding numerical study.
